% ============================================================
%  SEKCJA 1: WPROWADZENIE
% ============================================================

\section{Wprowadzenie}

Miód jest jednym z~najstarszych produktów spożywczych w~historii
ludzkości, a~jednocześnie jednym z~najczęściej fałszowanych surowców
na~współczesnym rynku żywności. Według danych Europejskiej Bazy Danych
Żywności (FoodFraudNetwork) miód zajmuje trzecie miejsce na~liście
produktów o~najwyższym ryzyku fałszowania, ustępując jedynie mleku
i~oliwie z~oliwek~\cite{apimondia_fraud_v2}. W~2023~roku łączny
wolumen importu miodu do~Unii Europejskiej wyniósł ponad 360~tys.~ton
o~wartości zbliżonej do~1~mld~EUR~\cite{cbi2024}, przy jednoczesnej
produkcji własnej UE wynoszącej około 286~tys.~ton, pokrywającej
zaledwie 60\%~krajowego zapotrzebowania~\cite{pollinatorhub2025}.
Strukturalny deficyt produkcyjny UE sprawia, że rynek jest w~znacznej
mierze uzależniony od~dostaw z~krajów trzecich, w~tym z~Chin (35\%
importu), Ukrainy (31\%) i~Argentyny (12\%).
Skala tego deficytu, połączona ze~znaczną dysproporcją cen między
miodem autentycznym a~substytutami syropowymi, tworzy  bodźce ekonomiczne do~podejmowania praktyk fałszerskich~\cite{apimondia_fraud_v2}.

W obliczu tych wyzwań kluczową rolę odgrywa Apimondia (Międzynarodowa Federacja Związków Pszczelarskich). Ta zrzeszająca ponad 80 krajowych organizacji organizacja z siedzibą w Rzymie stanowi globalne forum współpracy, którego celem jest promowanie naukowego, technicznego i ekonomicznego postępu w pszczelarstwie~\cite{apimondia_web}.

Stanowisko Apimondia opublikowane w~2025~roku przez Grupę Roboczą ds.~Fałszowania Produktów Pszczelich stanowi
istotny głos w~debacie o~integralności produktu~\cite{garcia2025}.
Autor identyfikuje systemową produkcję niedojrzałego miodu, definiowaną jako celowy odbiór miodu przed zakończeniem procesu
dojrzewania biologicznego w~ulu, połączony z~jego przemysłową dehydratacją, jako oddzielną i~poważną formę fałszowania, odrębną
od~klasycznych praktyk adulteracji syropami cukrowymi~\cite{garcia2025}.
Znaczenie tego stanowiska jest istotne: Apimondia, jako największa światowa federacja organizacji pszczelarskich,
dysponuje realnym wpływem na~kształtowanie zarówno norm międzynarodowych, jak i~praktyk kontrolnych stosowanych przez
inspekcje sanitarne w~krajach członkowskich UE.

Niniejszy artykuł nie kwestionuje istnienia ani powagi systemowej produkcji niedojrzałego miodu jako zjawiska naruszającego zasady
uczciwego obrotu żywnością. \textbf{Celem pracy jest natomiast wskazanie,}
że~stanowisko Apimondia dokonuje, być może mimowolnie, nadmiernego
uogólnienia, utożsamiając ze~sobą trzy jakościowo odmienne zjawiska:
(1)~systemową, celową produkcję niedojrzałego miodu dla zysku
ekonomicznego; (2)~wczesny odbiór miodu w~warunkach wymuszonych przez
czynniki klimatyczne lub technologiczne; oraz (3)~obecność drożdży
osmofilnych jako naturalnych składników miodu~\cite{garcia2025,munitis1976}.
Konsekwencją tego uogólnienia jest tworzenie kryterium dyskwalifikującego opartego na~obecności drożdży, co~pozostaje
w~sprzeczności zarówno z~obowiązującymi standardami regulacyjnymi, jak i~aktualnym stanem wiedzy mikrobiologicznej~\cite{codex_stan12, rodriguezmachado2024}.

Analiza obowiązujących standardów międzynarodowych, w~tym Codex
Alimentarius CODEX STAN 12-1981, Dyrektywy UE 2001/110/WE oraz
projektu normy ISO/DIS 24607:2025, prowadzi wniosku: żaden z~tych dokumentów nie~ustanawia limitu liczby komórek drożdży
(CFU/g) jako kryterium oceny jakości lub autentyczności miodu~\cite{codex_stan12}.
Jedynym prawnym kryterium mikrobiologicznym jest \emph{fermentacja jako stan aktywny}, nie zaś sama obecność
mikroorganizmów zdolnych do~fermentacji~\cite{codex_stan12}.
Tymczasem aktywność fermentacyjna drożdży osmofilnych, w~tym
gatunków \textit{Zygosaccharomyces rouxii} i~\textit{Z.~richteri},
endemicznie obecnych w~miodzie wszystkich typów botanicznych i~geograficznych jest biologicznie niemożliwa przy aktywności wody $a_w < 0{,}60$, co~odpowiada wilgotności miodu na~poziomie $\leq 20\%$~\cite{zamora2006}.

Powyższe przesłanki prowadzą do~trójczłonowej tezy niniejszego artykułu:
po~pierwsze, obecność drożdży osmofilnych w~miodzie spełniającym
kryterium wilgotności $\leq 20\%$ jest normatywnie neutralna, gdyż
aktywność wody $a_w < 0{,}60$ uniemożliwia wzrost i~fermentację tych
mikroorganizmów; po~drugie, żaden obowiązujący standard krajowy
ani~międzynarodowy nie~penalizuje samej obecności drożdży bez~aktywnej
fermentacji, a~zatem samo stwierdzenie ich liczby metodą płytkową (CFU/g)
nie~stanowi podstawy prawnej ani naukowej do~dyskwalifikacji produktu
z~obrotu; po~trzecie, koncepcja \emph{dojrzałości funkcjonalnej},
rozumianej jako zintegrowana stabilność mikrobiologiczna, chemiczna
i~biologiczna, stanowi precyzyjniejsze i~operacjonalizowalne narzędzie
oceny jakości miodu niż wyłączne kryterium morfologiczne zasklepienia
komórek plastra~\cite{zhang2021,guo2020,sun2021}.

Artykuł ma~charakter przeglądowy i~obejmuje następujące obszary:
(i)~analizę międzynarodowych ram regulacyjnych w~zakresie parametrów mikrobiologicznych miodu; 
(ii)~biologię i~metabolomikę procesu dojrzewania miodu jako continuum; 
(iii)~mikrobiologię drożdży osmofilnych w~kontekście aktywności wody i~ryzyka fermentacji; 
(iv)~zróżnicowanie globalnych systemów produkcji pszczelarskiej i~ich implikacje dla~standaryzacji; 
(v)~analizę rynku miodu UE i~konsekwencji regulacyjnych nadmiernie rygorystycznych interpretacji; oraz
(vi)~propozycję modelu dojrzałości funkcjonalnej wraz z~rekomendacjami dla~organów normalizacyjnych i~Apimondia.